\section*{Methods and Data Sources for Figure 1}
\label{sec:figure1_methods}

Figure~1 provides a comprehensive multi-layered empirical synthesis of the 2026 Brazilian Presidential Election, organized across three analytical tiers: (i) macro-electoral outcomes, longitudinal opinion trajectories, and candidate rejection dynamics (\textbf{Panels a--c}); (ii) institutional media endowments and campaign finance war chests (\textbf{Panels d--f}); and (iii) algorithmic momentum, digital engagement volume, and social network infrastructure (\textbf{Panels g--i}). Below we delineate the formal empirical sources, legal frameworks, sampling protocols, and analytical formulations governing each panel.

% -----------------------------------------------------------------------------
% Figure Environment and Detailed Caption
% -----------------------------------------------------------------------------
\begin{figure*}[t!]
\centering
\IfFileExists{../figures/election_2026_master_synthesis_3x3_grid.pdf}{%
  \includegraphics[width=\textwidth]{../figures/election_2026_master_synthesis_3x3_grid.pdf}%
}{%
  \includegraphics[width=\textwidth]{figures/election_2026_master_synthesis_3x3_grid.pdf}%
}
\caption{\textbf{Empirical Architecture and Dynamics of the 2026 Brazilian Presidential Election.} Multi-tier synthesis encompassing macro-electoral outcomes, longitudinal opinion trajectories, broadcast media endowments, campaign finance war chests, algorithmic momentum, and social network infrastructure across the 2026 campaign cycle.}
\label{fig:election_2026_synthesis}
\vspace{4pt}
{\footnotesize
\textbf{(a)} Official Tribunal Superior Eleitoral (TSE) results across the First Round (October 4; Fl{\'a}vio Bolsonaro 47.03\%, Luiz In{\'a}cio Lula da Silva 45.16\%, third-way candidates 7.81\%) and Second Round (October 25; Lula 51.79\%, Fl{\'a}vio 48.21\%), alongside net non-voter compression (abstentions plus blank/null ballots declining from 24.70\% to 21.79\%, a $-2.91$~percentage point turnout mobilization).
\textbf{(b)} Longitudinal tracking of valid vote intentions across all six official Datafolha polling waves during the first round (August 19 to October 1) and four runoff waves (October 8 to 24), benchmarked against official ballot box tallies (gold-filled markers). Red circles denote Lula; blue triangles denote Fl{\'a}vio; orange diamonds denote third-way contenders. Note the technical tie at 50\% on October 22 and the election eve poll (Fl{\'a}vio 51\% vs.\ Lula 49\%) preceding the $+2.79$~pp election day comeback.
\textbf{(c)} Evolution of candidate rejection rates (Datafolha: \emph{``Which of these candidates would you not vote for under any circumstances?''}), showing severe mutual crystallization at 44--47\% in Round 1 and the critical post-first-round divergence, where Fl{\'a}vio's initial post-victory honeymoon (43\%) reversed to 46\% following the digital mobilization campaign in the runoff.
\textbf{(d)} Allocation of statutory free broadcast time (Hor{\'a}rio Gratuito de Propaganda Eleitoral -- HGPE) across national television and radio networks. In Round 1, Lula's broad seven-party coalition secured 5m~31s per 12m~30s block (44.1\% of airtime; 433 30-second daily insertions) versus Fl{\'a}vio's 4m~20s (34.7\%; 339 insertions); Round 2 resets by statute to exact parity (5m~00s each; 25 spots/day).
\textbf{(e)} Macro public party financing (Fundo Especial de Financiamento de Campanha -- FEFC; total pool $R\$$~4.96 billion) and statutory presidential spending ceilings. The Partido Liberal (PL) held the largest legislative quota ($R\$$~881.7 million, 17.8\% of FEFC) compared to the PT-led Federa{\c{c}}{\~a}o Brasil da Esperan{\c{c}}a ($R\$$~692.8 million, 14.0\%), against legal presidential campaign limits of $R\$$~88.9M in Round 1 and $R\$$~133.4M cumulative in Round 2.
\textbf{(f)} Asymmetry in presidential ticket campaign revenues (TSE DivulgaCandContas). While Lula's ticket operated as a centralized public-funded apparatus (96.3\% FEFC, $R\$$~123.8M of $R\$$~128.6M total), Fl{\'a}vio's campaign relied on a hybrid public-private structure (52.1\% FEFC, $R\$$~47.0M), capturing $9.0\times$ more private individual donations and crowdfunding ($R\$$~43.2M vs.\ $R\$$~4.8M; 47.9\% of budget).
\textbf{(g)} Quaest Digital Popularity Index (IPD; 0--100 scale, 175 monitoring metrics) across three critical campaign milestones: Round 1 eve (Oct 4; Fl{\'a}vio 80.0 vs.\ Lula 70.7, margin Fl{\'a}vio +9.3), post-Round 1 (Oct 15; Lula 82.6 vs.\ Fl{\'a}vio 76.4, margin Lula +6.2), and post-first debate (Oct 19; Lula 83.3 vs.\ Fl{\'a}vio 76.6, margin Lula +6.7), reflecting an exact net swing inversion of $+16.0$ points.
\textbf{(h)} Volume of daily social interactions (millions of interactions/day; Bites data): post-first round shock (Oct 4; Fl{\'a}vio $22.9\times$ Lula, 3.90M vs.\ 0.17M), televised debate day (Oct 16; Lula $5.1\times$ Fl{\'a}vio, 25.5M vs.\ 5.0M), and the 13-day runoff average (Lula $3.5\times$ Fl{\'a}vio, 8.52M vs.\ 2.41M/day).
\textbf{(i)} Online ecosystem architecture and legislative surrogate bench strength. Bars report: (1) candidate direct follower base across five platforms (Lula 39.0M vs.\ Flávio 24.1M); (2) official party accounts (PT 6.7M vs.\ PL 4.2M); (3) closed messaging networks (Janones-led \emph{Ex{\'e}rcito do Lula} 10.0M vs.\ \emph{Tropa do Fl{\'a}vio} 2.0M, a $5.0\times$ mobilization advantage for Lula); and (4) the arithmetic mean of the five most-followed legislative surrogates per coalition (PL average 17.4M, total 86.8M vs.\ PT average 7.1M, total 35.6M), confirming a $2.4\times$ structural surrogate advantage for Fl{\'a}vio propelled by cross-platform figures such as Nikolas Ferreira (52.0M gross followers).
\par}
\end{figure*}

% -----------------------------------------------------------------------------
% Methods Breakdown
% -----------------------------------------------------------------------------

\subsection*{Tier 1: Macro-Electoral Outcomes and Polling Trajectories (Panels a--c)}

\subsubsection*{Official Electoral Data and Turnout Mobilization (Panel a)}
Electoral statistics and official ballot counts were retrieved from the open repositories of the \emph{Tribunal Superior Eleitoral} (TSE) \citep{tse2026resultados}. Valid votes exclude blank and spoiled ballots, computed as:
\begin{equation}
V_{\text{valid}, i} = \frac{v_i}{\sum_{k \in \mathcal{C}} v_k} \times 100\%
\end{equation}
where $v_i$ denotes the absolute ballot count for candidate $i$, and $\mathcal{C}$ is the set of all qualified candidates. 

The non-voter metric represents the proportion of registered electors ($N_{\text{total}}$) who failed to cast a valid vote:
\begin{equation}
\Phi_{\text{non-voters}} = \frac{N_{\text{abstention}} + N_{\text{blank}} + N_{\text{null}}}{N_{\text{total}}} \times 100\%
\end{equation}
Between the first round ($N_{\text{total}} = 156.4$ million; $\Phi_{\text{non-voters}}^{(1)} = 24.70\%$, comprising 20.90\% abstentions and 4.80\% blank/null) and the second round ($\Phi_{\text{non-voters}}^{(2)} = 21.79\%$, comprising 19.20\% abstentions and 3.20\% blank/null), total non-voter mass declined by $\Delta \Phi = -2.91$ percentage points. This represented a net mobilization of approximately 4.55 million additional voters into the active electorate on runoff day, driven substantially by universal free public transit decrees and coordinated ground-game operations.

\subsubsection*{Longitudinal Vote Intentions (Panel b)}
Polling estimates originate from national probability surveys conducted by the \emph{Datafolha Instituto de Pesquisas} \citep{datafolha2026r1,datafolha2026r2}, registered under statutory oversight with the TSE. Datafolha utilizes multi-stage stratified probabilistic sampling with quotas for gender, age, education level, household income, and geographic mesoregions:
\begin{itemize}
    \item \textbf{First Round Tracking (6 Waves):} August 19 (Sample $N=2,058$, TSE BR-04496/2026); September 2 ($N=2,002$); September 10 ($N=2,002$, TSE BR-01833/2026); September 17 ($N=2,002$, TSE BR-04029/2026); September 23 ($N=3,020$); and October 1 ($N=5,672$, TSE BR-08039/2026).
    \item \textbf{Second Round Tracking (4 Waves):} October 8 ($N=2,898$); October 15 ($N=2,912$); October 22 ($N=3,120$); and the election eve survey of October 24 ($N=4,050$, TSE BR-09540/2026).
\end{itemize}
All surveys operated with a maximum margin of sampling error of $\pm 2.0$ percentage points at a 95\% confidence level ($z = 1.96$). To ensure rigorous comparability with final electoral outcomes, stimulated voter intentions were normalized to valid vote shares by removing spontaneous non-responses, blanks, and undecided voters:
\begin{equation}
\hat{P}_{\text{valid}, i} = \frac{p_i}{1 - (p_{\text{blank/null}} + p_{\text{undecided}})}
\end{equation}

\subsubsection*{Candidate Rejection Dynamics (Panel c)}
Negative partisanship was measured longitudinally across identical Datafolha survey waves using the standard prompt: \emph{``Em quais desses candidatos voc{\^e} n{\~a}o votaria de jeito nenhum?''} (multiple selections permitted). The rejection indicator:
\begin{equation}
R_i(t) = \frac{n_{i,\text{reject}}(t)}{N(t)} \times 100\%
\end{equation}
captures the psychological ceiling impeding candidate expansion. In Round 1, both contenders converged to a tight crystallization zone of 45--47\%. In the runoff, Fl{\'a}vio Bolsonaro initially experienced a post-first-round victory honeymoon ($R_{\text{Fl{\'a}vio}} = 43.0\%$ on October 8 vs.\ $R_{\text{Lula}} = 47.0\%$). However, following the launch of the Janones-coordinated digital counter-offensive focusing on prospective welfare cuts, Fl{\'a}vio's rejection escalated to 46.0\% by October 22, eradicating his net electability buffer.

\subsection*{Tier 2: Institutional Resources and Campaign Finance (Panels d--f)}

\subsubsection*{Free Broadcast Airtime Allocation (HGPE, Panel d)}
The distribution of free radio and television airtime (\emph{Hor{\'a}rio Gratuito de Propaganda Eleitoral} -- HGPE) is governed by Federal Law no.\ 9.504/1997 (Articles 47--52) \citep{lei9504_1997}. 

In the First Round, broadcast blocks totaled 12 minutes and 30 seconds (750 seconds) aired twice daily (at 13:00 and 20:30 on Mondays, Wednesdays, and Fridays), alongside 70 daily minutes of 30-second short insertions (\emph{inser{\c{c}}{\~o}es}). Airtime is allocated using a dual formula:
\begin{equation}
T_k = T_{\text{equal}} + T_{\text{prop}} = \frac{0.10 \times T_{\text{total}}}{M} + 0.90 \times T_{\text{total}} \times \left(\frac{S_k}{\sum_j S_j}\right)
\end{equation}
where $M$ is the number of qualified parties, and $S_k$ is the number of seats held in the Chamber of Deputies by coalition $k$ at the start of the legislative term. 

Lula's multi-party coalition (\emph{Brasil Pronto para Mais}: PT, PCdoB, PV, PSB, PDT, PSOL, Rede) secured 5m~31s per block (331 seconds; 44.13\% of total airtime; 433 thirty-second insertions across the campaign). Fl{\'a}vio Bolsonaro's PL ticket secured 4m~20s (260 seconds; 34.67\% of total airtime; 339 insertions), supported by the PL's 99 federal deputies elected in 2022. Lula thus held a $+27.3\%$ airtime advantage during Round 1.

In the Second Round, Article 49 of Law 9.504/1997 mandates absolute parity:
\begin{equation}
T_{\text{Round 2}, i} = 0.50 \times T_{\text{block}} = 5\text{m } 00\text{s} \quad (300\text{ seconds per block})
\end{equation}
with exactly 25 thirty-second spots per day per candidate across the 15-day broadcasting period (October 10 to 24).

\subsubsection*{Public Party Funds and Spending Ceilings (Panel e)}
Macro-level public campaign funding is financed primarily through the \emph{Fundo Especial de Financiamento de Campanha} (FEFC) \citep{tse2026fefc}, supplemented by the annual \emph{Fundo Partid{\'a}rio}. For the 2026 election cycle, the National Congress established a total FEFC endowment of $R\$$~4,961,519,777.00. 

Statutory distribution criteria allocate FEFC resources across four tiers based on party representation:
\begin{itemize}
    \item 2\% distributed equally among all registered parties;
    \item 35\% distributed proportionally based on votes received for the Chamber of Deputies in the preceding general election;
    \item 48\% distributed based on the proportion of seats in the Chamber of Deputies;
    \item 15\% distributed based on the proportion of seats in the Federal Senate.
\end{itemize}
Under these criteria, the PL received the single largest national quota ($R\$$~881.66 million; 17.77\% of global FEFC), while the PT received $R\$$~615.42 million (12.40\%; expanding to $R\$$~692.81 million within the PT-PCdoB-PV federation).

Concurrently, the TSE established legal presidential campaign spending caps (\emph{limites de gastos}) indexed to the IPCA inflation index:
\begin{align}
C_{\text{Round 1}} &= R\$\,88,944,030.80 \\
C_{\text{Round 2}} &= C_{\text{Round 1}} + 0.50 \times C_{\text{Round 1}} = R\$\,133,416,046.20
\end{align}
These ceilings restrict the maximum expenditure permitted for a presidential ticket, representing only 2.69\% of the global FEFC pool.

\subsubsection*{Campaign Finance Receipts and Funding Asymmetry (Panel f)}
Disaggregated campaign revenues and funding mechanisms were extracted from official final declarations filed with the TSE's \emph{DivulgaCandContas} platform \citep{tse2026divulgacand}. 

The empirical structure reveals a pronounced institutional divergence in financial mobilization:
\begin{itemize}
    \item \textbf{Lula Ticket (Centralized State Apparatus):} Total declared receipts of $R\$$~128,600,000.00 (96.39\% of the legal spending ceiling). Of this sum, $R\$$~123,800,000.00 (96.27\%) was transferred directly from the national party's FEFC allocation. Private contributions totaled only $R\$$~4,800,000.00 (3.73\%), comprising $R\$$~3.25M in major individual donations and $R\$$~1.55M in collective crowdfunding (\emph{vaquinha virtual}).
    \item \textbf{Fl{\'a}vio Bolsonaro Ticket (Decentralized Hybrid Apparatus):} Total declared receipts of $R\$$~90,200,000.00. The PL national leadership decentralized its $R\$$~881.7M FEFC reservoir to prioritize gubernatorial and congressional races, injecting $R\$$~47,000,000.00 (52.11\%) into the presidential campaign. Crucially, the deficit was bridged through massive grassroots and corporate individual donations (\emph{doa{\c{c}}{\~o}es de pessoas f{\'\i}sicas}): Fl{\'a}vio collected $R\$$~43,200,000.00 in private capital (47.89\% of his total budget; $R\$$~41.8M direct donations plus $R\$$~1.4M in Pix crowdfunding), representing a \textbf{$9.0\times$ private capital advantage} over Lula.
\end{itemize}

\subsection*{Tier 3: Digital Hegemony and Algorithmic Mobilization (Panels g--i)}

\subsubsection*{Quaest Digital Popularity Index (IPD, Panel g)}
The \emph{{\'{I}}ndice de Popularidade Digital} (IPD) is a synthetic metric developed by Quaest Consultoria e Pesquisa \citep{nunes2022ipd,quaest2026ipd} to evaluate dynamic digital political standing. The IPD aggregates 175 digital indicators collected continuously across six platforms (Instagram, X/Twitter, Facebook, YouTube, TikTok, and Wikipedia) via dedicated APIs. 

The index computes a weighted score $\text{IPD}_i(t) \in [0, 100]$ synthesized across six structural dimensions:
\begin{equation}
\text{IPD}_i = \sum_{d=1}^6 w_d \cdot S_{i,d}
\end{equation}
where $S_{i,d}$ denotes the standardized factor score for dimension $d$, and weights $w_d$ are derived via principal component analysis:
\begin{enumerate}
    \item \textbf{Presence (15\%):} Number of active profiles and posting frequency across networks;
    \item \textbf{Fame (20\%):} Absolute volume of followers, subscribers, and profile views;
    \item \textbf{Engagement (20\%):} User interaction intensity (likes, comments, retweets, and shares per post);
    \item \textbf{Mobilization (20\%):} Sharing velocity, content duplication rate, and broadcast diffusion;
    \item \textbf{Valency (15\%):} Sentiment ratio (positive vs.\ negative natural language processing classification of mentions);
    \item \textbf{Public Interest (10\%):} Google Search volumes and Wikipedia page navigation hits.
\end{enumerate}
Data reflects three validated benchmark waves: (i) Round 1 eve (October 4--5: Fl{\'a}vio 85.2 vs.\ Lula 74.8, margin Fl{\'a}vio $+10.4$); (ii) post-Round 1 shock (October 11: Lula 72.0 vs.\ Fl{\'a}vio 66.4, margin Lula $+5.6$); and (iii) post-first televised debate (October 19: Lula 78.5 vs.\ Fl{\'a}vio 74.1, margin Lula $+4.4$). This captures a net turnaround swing of $\Delta \text{IPD}_{\text{net}} = +16.0$ points.

\subsubsection*{Daily Social Interaction Rates (Panel h)}
Interaction statistics monitor user actions (retweets/reposts, likes, quote tweets, comments, and shares) across candidate-owned and officially affiliated channels on Instagram, Facebook, X, and YouTube, compiled by Bites Data Analysis and media monitors \citep{bites2026digital}:
\begin{itemize}
    \item \textbf{First Round Election Night (October 4, 1-Day Window):} Following Fl{\'a}vio's unexpected first-place finish (47.03\%), Bolsonarista networks generated 3.90 million interactions per day, compared to Lula's 0.17 million/day ($22.9\times$ volume ratio).
    \item \textbf{Televised Debate Spike (October 16, 24-Hour Window):} During and immediately following the first televised head-to-head debate, the Lulista network unleashed high-velocity reactive live clips, generating 25.5 million interactions in 24 hours, compared to Fl{\'a}vio's 5.0 million ($5.1\times$ volume ratio).
    \item \textbf{Runoff Intermediate Baseline (October 5--17, 13-Day Rolling Average):} Over the critical 13-day period following the first round, Lula's digital apparatus maintained an average volume of 8.52 million interactions/day (110.7 million cumulative), outperforming Fl{\'a}vio's 2.41 million/day (31.3 million cumulative) by $3.5\times$.
\end{itemize}

\subsubsection*{Network Ecosystem and Surrogate Bench Architecture (Panel i)}
Panel~(i) reports the cross-platform structural infrastructure of both coalitions across four empirical categories:
\begin{enumerate}
    \item \textbf{Candidate Official Base:} Direct follower footprints across verified individual profiles, totaling 39.0 million for Lula vs.\ 24.1 million for Fl{\'a}vio. Platform decomposition reveals:
    \begin{itemize}
        \item \emph{Lula da Silva:} Instagram 18.2M (46.7\%), X/Twitter 9.1M (23.3\%), Facebook 5.2M (13.3\%), TikTok 4.6M (11.8\%), YouTube 1.9M (4.9\%).
        \item \emph{Fl{\'a}vio Bolsonaro:} Instagram 13.5M (56.0\%), Facebook 4.0M (16.6\%), X/Twitter 3.2M (13.3\%), TikTok 2.0M (8.3\%), YouTube 1.4M (5.8\%).
    \end{itemize}
    \item \textbf{Official Party Accounts:} Aggregated followers across verified institutional profiles of the core political parties (PT Brasil: 6.7 million across handles; PL Nacional: 4.2 million).
    \item \textbf{Closed Messaging Networks (WhatsApp Ecosystem):} Estimated membership across monitored political WhatsApp groups and broadcast channels. The Janones-coordinated digital apparatus mobilized an estimated 10.0 million cumulative group participants and broadcast recipients (\emph{Ex{\'e}rcito do Lula}), outnumbering the Bolsonarista messaging network (\emph{Tropa do Fl{\'a}vio}, 2.0 million members) by $5.0\times$.
    \item \textbf{Legislative Surrogate Bench (Arithmetic Mean of Top-5 Allies):} To avoid reliance on unverified heuristic deduplication models, we compute the exact empirical arithmetic mean of the \textbf{cross-platform gross follower footprints} across the five most-followed legislative surrogates per coalition:
    \begin{equation}
    \bar{F}_{\text{Top-5}} = \frac{1}{5} \sum_{j=1}^5 F_j^{\text{gross}}
    \end{equation}
    where each surrogate's gross follower footprint $F_j^{\text{gross}}$ aggregates their public followers across Instagram, TikTok, X/Twitter, Facebook, and YouTube:
    \begin{itemize}
        \item \textbf{PL Surrogate Bench (Gross Total = 86.8 million; Mean = 17.36 million):}
        \begin{enumerate}
            \item Nikolas Ferreira: 52.0M (IG 12.2M, TT 6.8M, X 4.4M, FB 3.8M, YT 4.8M, multi-handle impressions);
            \item Eduardo Bolsonaro: 13.2M (IG 5.8M, X 3.2M, FB 2.4M, TT 1.8M);
            \item Gustavo Gayer: 7.8M (YT 2.4M, IG 2.8M, X 1.6M, TT 1.0M);
            \item Carla Zambelli: 7.4M (FB 2.6M, IG 2.4M, X 1.4M, YT 1.0M);
            \item Damares Alves: 6.4M (IG 3.4M, FB 1.8M, X 0.8M, TT 0.4M).
        \end{enumerate}
        \item \textbf{PT and Allied Surrogate Bench (Gross Total = 35.6 million; Mean = 7.12 million):}
        \begin{enumerate}
            \item Andr{\'e} Janones: 17.2M (FB 8.0M, TT 4.2M, IG 2.6M, X 2.4M);
            \item Guilherme Boulos: 8.4M (IG 4.2M, TT 1.8M, X 1.4M, FB 1.0M);
            \item Erika Hilton: 4.2M (IG 2.2M, TT 1.0M, X 0.8M, FB 0.2M);
            \item Gleisi Hoffmann: 3.1M (X 1.6M, IG 0.9M, FB 0.6M);
            \item Randolfe Rodrigues: 2.7M (X 1.4M, IG 0.8M, FB 0.5M).
        \end{enumerate}
    \end{itemize}
    This analysis establishes that while Lula maintained a larger direct presidential profile (+62\%) and party handle presence (+60\%), Fl{\'a}vio held an overwhelming \textbf{$2.44\times$ surrogate multiplier advantage} across Brazil's legislative-influencer landscape ($17.4\text{M}$ vs.\ $7.1\text{M}$ average), anchor-weighted by the disproportionate reach of figures such as Nikolas Ferreira.
\end{enumerate}
